\documentclass[conference]{IEEEtran}
\IEEEoverridecommandlockouts
% The preceding line is only needed to identify funding in the first footnote. If that is unneeded, please comment it out.
\usepackage{cite}
\usepackage{amsmath,amssymb,amsfonts}
\usepackage{algorithmic}
\usepackage{graphicx}
\usepackage{textcomp}
\usepackage{xcolor}
\usepackage{nameref}
\usepackage{autoref}
\usepackage{hyperref}
\def\BibTeX{{\rm B\kern-.05em{\sc i\kern-.025em b}\kern-.08em
    T\kern-.1667em\lower.7ex\hbox{E}\kern-.125emX}}
\newcommand*{\fullref}[1]{\textit{\hyperref[{#1}]{\autoref*{#1} \nameref*{#1}}}}

\begin{document}

\title{Duplicate recognition for restaurant dataset*\\
{\footnotesize \textsuperscript{*}Note: Sub-titles are not captured in Xplore and
should not be used}
}

\author{\IEEEauthorblockN{1\textsuperscript{st} Alexander Christoph}
\IEEEauthorblockA{\textit{OTH Regensburg} \\
Regensburg, Germany \\
alexander.christoph@st.oth-regensburg.de}
}

\maketitle

\begin{abstract}
This document describes the analysis and removal of duplicates from the restaurants dataset. The aim is to remove as many duplicates as possible from the dataset and store the data without duplicates in a cloud hosted mongodb instance. A problem when finding duplicates of restaurants (or nearly any other dataset) is the format and the different writing of the entries in the data. This problem was already researched by several IEEE members (quelle). Within my research there were made different approaches to remove the duplicates which are described below. The accuracy of the results is measured with precision, recall and F-score. After removing the duplicates the cleared dataset is stored into a mongodb cluster so that it can be accessed any time. In this paper I will also describe some techniques which I haven't used in my project but are also very useful. 

\end{abstract}

\begin{IEEEkeywords}
component, formatting, style, styling, insert
\end{IEEEkeywords}

\section{Introduction}
In times where big data gets more and more attraction from the industry it is very important to learn how to deal with it. Especially when it comes to the structure and format of the data. When looking at big data, it is most of the time a problem that there are duplicates and unclean entries in a dataset. This gives inaccurate results when analyzing or working with this data. That happens because most of the time there isn't that much preprocessing happening and there aren't even checks for a standard data format. To learn how to deal with duplicate data, the restaurants dataset is used. This isn't actually big data, but to understand the importance of the preprocessing task it is pretty good because it's considered as a well researched dataset to play with and compare to the gold standard. 

In my research I've looked into different approaches to detect duplicates and remove them from the dataset. The first approach was to just remove all duplicates that are in the data, this wasn't successful because most of the duplicates have different writings or completely different values in some of the fields. So I started to analyze the data and look for potential duplicates and how to prepare them so that the program can match them. The first step was to remove all special characters and some other unnecessary contents in the different columns. After that I investigated which columns are the most useful when it comes to duplicate detection. 

After researching and removing potential duplicates, I calculated the count for true positives, false positives, true negatives and false negatives of my prediction with the help of the gold standard duplicate dataset which was evaluated by hand. With the help of these metrics I calculated the recall and precision for my result to get a better understanding, how good my evaluations were. As a conclusion the values I got were: 
\begin{itemize}
	\item All entries in original dataset: 864
	\item Detected duplicates (all): 111
	\item Real duplicates (from gold standard): 112
	\item True positives: 103
	\item True negatives: 744
	\item False positives: 8
	\item False negatives: 9
	\item Precision: 0.93
	\item Recall: 0.92
\end{itemize}

After the methods are applied and the duplicates are removed it is necessary to store the new dataset somewhere. For this I have choosen mongodb because of it's great compatibility with many programming languages and the low expenses when you want to store data in it. Mongodb could also be used for many preprocessing tasks because of the great aggregation framework that it offers. 
\section{The restaurants dataset}
The restaurants dataset which is researhed in this paper is a .tsv dataset which contains 864 rows of data with six columns. The columns of the dataset are: 
\begin{itemize}
	\item id: The unique id of each row
	\item name: The name of the restaurant
	\item address: The address where the restaurant is located
	\item city: The city of the restaurant
	\item phone: The phone number of the restaurant
	\item type: The kind of the restaurant (i.e. french or american)
\end{itemize}
In the data there are 122 duplicated restaurants. These duplicates were picked by hand from some researchers to define a gold standard which would be the best possible result after removing all duplicates. The duplicates that occur in the dataset have different deviations from each other. For example some duplicates have a different order of the words in their name field like "the palm" and "palm the". Others have different separators for the phone number like "310/659-9639" and "310-659-9639". Sometimes the city field of the duplicates is a district from a bigger city and sometimes it's the city name itself like "los angeles" and "hollywood". There are even more  different deviations in the dataset as well whose solution to detect them will be discussed later in \fullref{sec_methods} section.

In my research I focused mostly on the columns name, city, address and phone because they have the most useful information when it comes to duplicate detection.The id column was left out because it's only a unique identifier which wouldn't bring a benefit for duplicate recognition. The type column was left out because there are to much restaurants with the same type which would result in an unclean target data record.

Besides of the plain restaurants dataset there is also a dataset given which contains all the duplicates in the data by id. In this duplicate set, there are only two columns, "id1" and "id2" which define the original id and the duplicate id. A dataset without all these duplicates is considered as gold standard. The reached results will be measured to this gold standard.
\section{Metrics used to measure th results}
To measure the accuracy of results when it comes to duplicate detection there are different metrics that are commonly used. For all these metrics it's important to pre calculate four values with the help of the gold standard data and the archived results: 
\begin{itemize}
	\item True positives (TP): The correct classified true entries
	\item True negatives (TN): The correct classified false entries
	\item False positives (FP): Incorrect classified true entries
	\item False negatives (FN): Incorrect classified false entries
\end{itemize}
In a gold standard dataset there are no false positives and false negatives.

A first metric that is important is the well known accuracy. 
\begin{align}
 accuracy =  \frac{TP + TN}{TP + TN + FP + FN}
\end{align}
Only calculating the accuracy isn't enough in this task because the it only gives reliable results when a dataset is balanced. In this case it means that there are as many duplicated entries as non duplicates. As an addition for the accuracy I use two other metrics, precision and recall.

Precision measures the exactness for the minority class by only considering the positive classified entries of the result set. Because of this it is a good measurement for unbalanced data.
\begin{align}
	Precision = \frac{TruePositives}{TruePositives + FalsePositives}
\end{align}
Recal instead gives the accuracy for the fraction of relevant data.
\begin{align}
	Recall = \frac{TruePositives}{TruePositives + FalseNegatives}
\end{align}
\section{Methods used to detect duplicates} \label{sec_methods}
The research was written in python, a programming language, which is widely used in data science tasks. For my research I mostly used the Python pandas library, which has good data analysis features. 

As a whole there are 112 duplicates in the dataset which were selected by hand.

At first I only looked into the first 100 entries and tried to understand the difference and connections between the different rows and columns. I found out that many noise in the data came from special characters. So I removed every special character except spaces from the four columns that I looked into. Large and lower case letters didn't have to be considered because the whole data was in lower case letters. The received metrics after the first step of the data cleaning pipeline were: 
\begin{itemize}
	\item True positives: 80
	\item True negatives: 746
    \item False positives: 6
	\item False negatives: 32
	Accuracy 0.96
    \item Precision: 0.9302325581395349
	\item Recall: 0.7142857142857143
\end{itemize}
This is quite good for simply removing special characters, but the results can be optimized more. 

The next step of my data cleaning pipeline was to map multiple occurrences of the same city with different writing to as single key. For example there were entries which had "la" as a value and others had "los angeles". After applying this, the metrics were: 
\begin{itemize}
	\item Detected duplicates: 106
	\item True positives: 98
	\item True negatives: 746
	\item False positives: 6
	\item False negatives: 14
	\item Accuracy: 0.98
	\item Precision: 0.94
	\item Recall: 0.88
\end{itemize}
It's noticeable that the results have improved slightly after mapping the citynames. As a next step I removed appendixes from the address column. For this I have chosen to remove everything which occurs after "between", "off", "near", "at" or "in" from thee address string because many addresses had more precise descriptions for the address after the real address. Following up, I decided to remove appendixes of the street number like "1st" to "1" or "2nd" to 2. Then I looked again on the dataset and found out that the columns "address", as well as "name" both sometimes have a direction (like "north" or in short "n") added and sometimes don't, which is very inconsistent. So I removed every occurrence of a direction from this two columns. Neither of these three actions affected the metrics. 

So I researched the dataset again and found out that there are more inconsistencies in the address field. At first I discovered that some numbers in the address were written ass tring and others as numbers. So I mapped these to only represent numbers (i.e. "first" to 1). I also noticed that some words sometimes were written in short in the address column. These were: 
\begin{itemize}
	\item los angeles: la
	\item avenue: ave
	\item road: rd
	\item boulevard: blv, blvd
	\item street: st
\end{itemize}
After remapping these, the metrics had a big improvement to: 
\begin{itemize}
	\item Detected duplicates: 111
	\item True positives: 103
	\item True negatives: 744
	\item False positives: 8
	\item False negatives: 9
	\item Precision: 0.93
	\item Recall: 0.92
\end{itemize}
These were the final metrics for my research. There are even more ways to improve the duplicate recognition, which will be discussed in the next section \fullref{sec_improvements}
\section{Potential improvements and other researches} \label{sec_imporvements}
\section{Conclusion}


\newpage 
1
\newpage
2
\newpage 
3
\newpage 
4
\newpage 
5
\newpage 
\newpage 
\newpage 
\section{Introduction}
This document is a model and instructions for \LaTeX.
Please observe the conference page limits. 

\section{Ease of Use}

\subsection{Maintaining the Integrity of the Specifications}

The IEEEtran class file is used to format your paper and style the text. All margins, 
column widths, line spaces, and text fonts are prescribed; please do not 
alter them. You may note peculiarities. For example, the head margin
measures proportionately more than is customary. This measurement 
and others are deliberate, using specifications that anticipate your paper 
as one part of the entire proceedings, and not as an independent document. 
Please do not revise any of the current designations.

\section{Prepare Your Paper Before Styling}
Before you begin to format your paper, first write and save the content as a 
separate text file. Complete all content and organizational editing before 
formatting. Please note sections \ref{AA}--\ref{SCM} below for more information on 
proofreading, spelling and grammar.

Keep your text and graphic files separate until after the text has been 
formatted and styled. Do not number text heads---{\LaTeX} will do that 
for you.

\subsection{Abbreviations and Acronyms}\label{AA}
Define abbreviations and acronyms the first time they are used in the text, 
even after they have been defined in the abstract. Abbreviations such as 
IEEE, SI, MKS, CGS, ac, dc, and rms do not have to be defined. Do not use 
abbreviations in the title or heads unless they are unavoidable.

\subsection{Units}
\begin{itemize}
\item Use either SI (MKS) or CGS as primary units. (SI units are encouraged.) English units may be used as secondary units (in parentheses). An exception would be the use of English units as identifiers in trade, such as ``3.5-inch disk drive''.
\item Avoid combining SI and CGS units, such as current in amperes and magnetic field in oersteds. This often leads to confusion because equations do not balance dimensionally. If you must use mixed units, clearly state the units for each quantity that you use in an equation.
\item Do not mix complete spellings and abbreviations of units: ``Wb/m\textsuperscript{2}'' or ``webers per square meter'', not ``webers/m\textsuperscript{2}''. Spell out units when they appear in text: ``. . . a few henries'', not ``. . . a few H''.
\item Use a zero before decimal points: ``0.25'', not ``.25''. Use ``cm\textsuperscript{3}'', not ``cc''.)
\end{itemize}

\subsection{Equations}
Number equations consecutively. To make your 
equations more compact, you may use the solidus (~/~), the exp function, or 
appropriate exponents. Italicize Roman symbols for quantities and variables, 
but not Greek symbols. Use a long dash rather than a hyphen for a minus 
sign. Punctuate equations with commas or periods when they are part of a 
sentence, as in:
\begin{equation}
a+b=\gamma\label{eq}
\end{equation}

Be sure that the 
symbols in your equation have been defined before or immediately following 
the equation. Use ``\eqref{eq}'', not ``Eq.~\eqref{eq}'' or ``equation \eqref{eq}'', except at 
the beginning of a sentence: ``Equation \eqref{eq} is . . .''

\subsection{\LaTeX-Specific Advice}

Please use ``soft'' (e.g., \verb|\eqref{Eq}|) cross references instead
of ``hard'' references (e.g., \verb|(1)|). That will make it possible
to combine sections, add equations, or change the order of figures or
citations without having to go through the file line by line.

Please don't use the \verb|{eqnarray}| equation environment. Use
\verb|{align}| or \verb|{IEEEeqnarray}| instead. The \verb|{eqnarray}|
environment leaves unsightly spaces around relation symbols.

Please note that the \verb|{subequations}| environment in {\LaTeX}
will increment the main equation counter even when there are no
equation numbers displayed. If you forget that, you might write an
article in which the equation numbers skip from (17) to (20), causing
the copy editors to wonder if you've discovered a new method of
counting.

{\BibTeX} does not work by magic. It doesn't get the bibliographic
data from thin air but from .bib files. If you use {\BibTeX} to produce a
bibliography you must send the .bib files. 

{\LaTeX} can't read your mind. If you assign the same label to a
subsubsection and a table, you might find that Table I has been cross
referenced as Table IV-B3. 

{\LaTeX} does not have precognitive abilities. If you put a
\verb|\label| command before the command that updates the counter it's
supposed to be using, the label will pick up the last counter to be
cross referenced instead. In particular, a \verb|\label| command
should not go before the caption of a figure or a table.

Do not use \verb|\nonumber| inside the \verb|{array}| environment. It
will not stop equation numbers inside \verb|{array}| (there won't be
any anyway) and it might stop a wanted equation number in the
surrounding equation.

\subsection{Some Common Mistakes}\label{SCM}
\begin{itemize}
\item The word ``data'' is plural, not singular.
\item The subscript for the permeability of vacuum $\mu_{0}$, and other common scientific constants, is zero with subscript formatting, not a lowercase letter ``o''.
\item In American English, commas, semicolons, periods, question and exclamation marks are located within quotation marks only when a complete thought or name is cited, such as a title or full quotation. When quotation marks are used, instead of a bold or italic typeface, to highlight a word or phrase, punctuation should appear outside of the quotation marks. A parenthetical phrase or statement at the end of a sentence is punctuated outside of the closing parenthesis (like this). (A parenthetical sentence is punctuated within the parentheses.)
\item A graph within a graph is an ``inset'', not an ``insert''. The word alternatively is preferred to the word ``alternately'' (unless you really mean something that alternates).
\item Do not use the word ``essentially'' to mean ``approximately'' or ``effectively''.
\item In your paper title, if the words ``that uses'' can accurately replace the word ``using'', capitalize the ``u''; if not, keep using lower-cased.
\item Be aware of the different meanings of the homophones ``affect'' and ``effect'', ``complement'' and ``compliment'', ``discreet'' and ``discrete'', ``principal'' and ``principle''.
\item Do not confuse ``imply'' and ``infer''.
\item The prefix ``non'' is not a word; it should be joined to the word it modifies, usually without a hyphen.
\item There is no period after the ``et'' in the Latin abbreviation ``et al.''.
\item The abbreviation ``i.e.'' means ``that is'', and the abbreviation ``e.g.'' means ``for example''.
\end{itemize}
An excellent style manual for science writers is \cite{b7}.

\subsection{Authors and Affiliations}
\textbf{The class file is designed for, but not limited to, six authors.} A 
minimum of one author is required for all conference articles. Author names 
should be listed starting from left to right and then moving down to the 
next line. This is the author sequence that will be used in future citations 
and by indexing services. Names should not be listed in columns nor group by 
affiliation. Please keep your affiliations as succinct as possible (for 
example, do not differentiate among departments of the same organization).

\subsection{Identify the Headings}
Headings, or heads, are organizational devices that guide the reader through 
your paper. There are two types: component heads and text heads.

Component heads identify the different components of your paper and are not 
topically subordinate to each other. Examples include Acknowledgments and 
References and, for these, the correct style to use is ``Heading 5''. Use 
``figure caption'' for your Figure captions, and ``table head'' for your 
table title. Run-in heads, such as ``Abstract'', will require you to apply a 
style (in this case, italic) in addition to the style provided by the drop 
down menu to differentiate the head from the text.

Text heads organize the topics on a relational, hierarchical basis. For 
example, the paper title is the primary text head because all subsequent 
material relates and elaborates on this one topic. If there are two or more 
sub-topics, the next level head (uppercase Roman numerals) should be used 
and, conversely, if there are not at least two sub-topics, then no subheads 
should be introduced.

\subsection{Figures and Tables}
\paragraph{Positioning Figures and Tables} Place figures and tables at the top and 
bottom of columns. Avoid placing them in the middle of columns. Large 
figures and tables may span across both columns. Figure captions should be 
below the figures; table heads should appear above the tables. Insert 
figures and tables after they are cited in the text. Use the abbreviation 
``Fig.~\ref{fig}'', even at the beginning of a sentence.

\begin{table}[htbp]
\caption{Table Type Styles}
\begin{center}
\begin{tabular}{|c|c|c|c|}
\hline
\textbf{Table}&\multicolumn{3}{|c|}{\textbf{Table Column Head}} \\
\cline{2-4} 
\textbf{Head} & \textbf{\textit{Table column subhead}}& \textbf{\textit{Subhead}}& \textbf{\textit{Subhead}} \\
\hline
copy& More table copy$^{\mathrm{a}}$& &  \\
\hline
\multicolumn{4}{l}{$^{\mathrm{a}}$Sample of a Table footnote.}
\end{tabular}
\label{tab1}
\end{center}
\end{table}

\begin{figure}[htbp]
\centerline{\includegraphics{fig1.png}}
\caption{Example of a figure caption.}
\label{fig}
\end{figure}

Figure Labels: Use 8 point Times New Roman for Figure labels. Use words 
rather than symbols or abbreviations when writing Figure axis labels to 
avoid confusing the reader. As an example, write the quantity 
``Magnetization'', or ``Magnetization, M'', not just ``M''. If including 
units in the label, present them within parentheses. Do not label axes only 
with units. In the example, write ``Magnetization (A/m)'' or ``Magnetization 
\{A[m(1)]\}'', not just ``A/m''. Do not label axes with a ratio of 
quantities and units. For example, write ``Temperature (K)'', not 
``Temperature/K''.

\section*{Acknowledgment}

The preferred spelling of the word ``acknowledgment'' in America is without 
an ``e'' after the ``g''. Avoid the stilted expression ``one of us (R. B. 
G.) thanks $\ldots$''. Instead, try ``R. B. G. thanks$\ldots$''. Put sponsor 
acknowledgments in the unnumbered footnote on the first page.

\section*{References}

Please number citations consecutively within brackets \cite{b1}. The 
sentence punctuation follows the bracket \cite{b2}. Refer simply to the reference 
number, as in \cite{b3}---do not use ``Ref. \cite{b3}'' or ``reference \cite{b3}'' except at 
the beginning of a sentence: ``Reference \cite{b3} was the first $\ldots$''

Number footnotes separately in superscripts. Place the actual footnote at 
the bottom of the column in which it was cited. Do not put footnotes in the 
abstract or reference list. Use letters for table footnotes.

Unless there are six authors or more give all authors' names; do not use 
``et al.''. Papers that have not been published, even if they have been 
submitted for publication, should be cited as ``unpublished'' \cite{b4}. Papers 
that have been accepted for publication should be cited as ``in press'' \cite{b5}. 
Capitalize only the first word in a paper title, except for proper nouns and 
element symbols.

For papers published in translation journals, please give the English 
citation first, followed by the original foreign-language citation \cite{b6}.

\begin{thebibliography}{00}
\bibitem{b1} G. Eason, B. Noble, and I. N. Sneddon, ``On certain integrals of Lipschitz-Hankel type involving products of Bessel functions,'' Phil. Trans. Roy. Soc. London, vol. A247, pp. 529--551, April 1955.
\bibitem{b2} J. Clerk Maxwell, A Treatise on Electricity and Magnetism, 3rd ed., vol. 2. Oxford: Clarendon, 1892, pp.68--73.
\bibitem{b3} I. S. Jacobs and C. P. Bean, ``Fine particles, thin films and exchange anisotropy,'' in Magnetism, vol. III, G. T. Rado and H. Suhl, Eds. New York: Academic, 1963, pp. 271--350.
\bibitem{b4} K. Elissa, ``Title of paper if known,'' unpublished.
\bibitem{b5} R. Nicole, ``Title of paper with only first word capitalized,'' J. Name Stand. Abbrev., in press.
\bibitem{b6} Y. Yorozu, M. Hirano, K. Oka, and Y. Tagawa, ``Electron spectroscopy studies on magneto-optical media and plastic substrate interface,'' IEEE Transl. J. Magn. Japan, vol. 2, pp. 740--741, August 1987 [Digests 9th Annual Conf. Magnetics Japan, p. 301, 1982].
\bibitem{b7} M. Young, The Technical Writer's Handbook. Mill Valley, CA: University Science, 1989.
\end{thebibliography}
\vspace{12pt}
\color{red}
IEEE conference templates contain guidance text for composing and formatting conference papers. Please ensure that all template text is removed from your conference paper prior to submission to the conference. Failure to remove the template text from your paper may result in your paper not being published.

\end{document}
